\chapter{Further Tree Structures} \label{app:furtherTreeStructures}

\section*{Container-Driven Tree} \label{app:containerdrivenTree}
The concept of Containers, as described in \Cref{sec:model:containers} and illustrated in \Cref{fig:tikz:PKG:Containers}, organises Nodes into groups and thus creates a set of trees, whereas the root relates to the Container and the leaves to the Nodes. The only inner Member\footnote{inner Member: Analogous to a classic \glqq{}inner node\grqq{} but called \glqq{}inner Member\grqq{} to avoid confusion.} of each tree refers to its root, representing the Container itself. As visualised in \Cref{fig:tikz:mathModel:containers}, the set of trees is joined by employing a passive Member as virtual organising entity and thus passive Member that connects to each root of the subtree to create a superior tree. It is imperativ to mention that not every Node has to be part of a Container and thus not every Node has to be part of this tree.
\begin{figure}[h]
    \centering
    \begin{tikzpicture}

    \draw [decorate,decoration={brace,amplitude=5pt,raise=2ex}]
        (-6,2.6) -- ++(0,1.7) node[midway,xshift=-3.5em]{subtrees} coordinate (tmp);
    % \draw [fill = tubaf-accent-13!90, opacity=.2] (-7.3,2.6) rectangle ++(14.6,1.7);

    \draw [decorate,decoration={brace,amplitude=5pt,raise=2ex}]
        (tmp) -- ++(0,1) node[midway,xshift=-4.7em]{virtual to join};
    % \draw [fill = tubaf-accent-13!90, opacity=.25] (tmp) ++(-.3,0) rectangle ++(14.6,1);


    \node [MP] (root)               at (0,5) {virtual root};

    \node [MP] (container1)         at (-3,4) {container1};
    \node [MA, G0] (controller1)    at (-4.75,3) {\phantom{controller}};

    \node [MA, G0] (turtlesim1)     at (-1.3,3) {\phantom{turtlesim}};
    \node [MP] (container2)            at (3,4) {container2};
    \node [MA, G0] (controller2)    at (1.25,3) {\phantom{controller}};
    \node [MA, G0] (turtlesim2)     at (4.7,3) {\phantom{turtlesim}};

    \draw [TreeContainer, ->] (root) -| (container2);
    \draw [TreeContainer, ->] (root) -| (container1);

    \draw [TreeContainer, ->] (container1) -| (controller1);
    \draw [TreeContainer, ->] (container1) -| (turtlesim1);

    \draw [TreeContainer, ->] (container2) -| (controller2);
    \draw [TreeContainer, ->] (container2) -| (turtlesim2);

\end{tikzpicture}
    \caption[Exemplary and Joined Container-Driven Tree]{Exemplary and Joined Container-Driven Tree \legend} \label{fig:tikz:mathModel:containers}
\end{figure}

\section*{Process-Driven Tree} \label{app:processdrivenTree}
Based on the model that represents processes and their structure is described in \Cref{sec:model:processes} and illustrated in \Cref{fig:tikz:PKG:Processes}. The process-driven tree is formed by traversing the tree structure of the processes in the opposite direction to their creation up to the initial direct start or launch (via a launch script described in \Cref{sec:model:processes}). This creates a set of trees consisting of active Members only. To create a superior tree, the roots of all trees of the set are joined by a virtual organisational entity and thus passive Member. Each Node is assigned to a single process and thus part of this tree, but a process may encapsulate multiple Nodes (concept of Node Composition, see \Cref{sec:model:containers}). Furthermore, as each Node is enabled to start further processes, those processes are also representable within the created tree. \Cref{fig:tikz:mathModel:processes} illustrates an exemplary process tree. Within the example, all Nodes where started by a single launch script.
Through the Utilisation of passive Members as organisational entities, the mathematical meta-model enables the creation of a superior tree that might even spread across multiple operating systems. E.g. a ROS2 system that is spread across multiple computers.

\begin{figure}[h]
    \centering
    \begin{tikzpicture}

    \draw [decorate,decoration={brace,amplitude=5pt,raise=2ex}]
        (-6,2.6) -- ++(0,1.8) node[midway,xshift=-3.8em]{subtrees} coordinate (tmp);

    \draw [decorate,decoration={brace,amplitude=5pt,raise=2ex}]
        (tmp) -- ++(0,1) node[midway,xshift=-4.7em]{virtual to join};

    \node [MP] (root)               at (0,5.25) {virtual root};

    \node [MA] (launchScript)       at (0,4) {launch file};

    \node [MA, G0] (controller1)    at (-4.5,3) {\phantom{controller}};
    \node [MA, G0] (turtlesim1)     at (-1.9,3) {\phantom{turtlesim}};

    \node [MA, G0] (controller2)    at (1.9,3) {\phantom{controller}};
    \node [MA, G0] (turtlesim2)     at (4.5,3) {\phantom{turtlesim}};

    \draw [TreeProcess, ->] (root) -- (launchScript);

    \draw [TreeProcess, ->] (launchScript) -| (controller1);
    \draw [TreeProcess, ->] (launchScript) -| (turtlesim1);

    \draw [TreeProcess, ->] (launchScript) -| (controller2);
    \draw [TreeProcess, ->] (launchScript) -| (turtlesim2);

\end{tikzpicture}
    \caption[Exemplary Process-Driven Tree]{Exemplary Process-Driven Tree (of the ROS2 System) \legend} \label{fig:tikz:mathModel:processes}
\end{figure}


\chapter{Further Setup Information}
\begin{table}[h]
    \centering
    \caption[Tracepoints Employed for the Proof-of-Concept]{Tracepoints Employed for the Proof-of-Concept \cite{github:ros2_tracing}} \label{tab:tracepoints}
    \begin{tblr}{lX}
        \toprule
        Tracepoint & Purpose \\
        \midrule
        ros2:rcl\_node\_init                        & Indicates the instantiation of a ROS2 Node.                                                   \\
        ros2:rcl\_publisher\_init                   & Indicates the instantiation of a ROS2 publisher.                                              \\
        ros2:rcl\_subscription\_init                & Indicates the instantiation of a ROS2 subscriber.                                             \\
        ros2:rcl\_service\_init                     & Indicates the instantiation of a ROS2 Service server.                                         \\
        ros2:rcl\_client\_init                      & Indicates the instantiation of a ROS2 Service client.                                         \\
        ros2:rcl\_timer\_init                       & Indicates the instantiation of a ROS2 Timer.                                                  \\
        ros2:rclcpp\_timer\_link\_node              & Links an initialised Timer to its respective Node.                                            \\
        ros2:rcl\_lifecycle\_state\_machine\_init   & Indicates the instantiation of a ROS2 state machine that is necessitated by managed Nodes.    \\
        ros2:rcl\_lifecycle\_transition             & Indicates a lifecycle transition of managed Nodes.                                          \\
        ros2:rcl\_publish                           & Indicates the publishing of a message.                                                        \\
        \bottomrule
    \end{tblr}
\end{table}



\chapter{Performance Test: Detailed Results}

\begin{figure}[h]
    \centering
    \begin{tikzpicture}
    \begin{axis}[
        xlabel={Frequency [Hz]},
        ylabel={Number of Publisher-Subscriber Pairs (n)},
        colormap name=colormap-latencies,
        colorbar,
        colorbar sampled,
        point meta min=0,
        point meta max=0.05,
        colorbar style={
            samples=512,
            ylabel=Mann-Whitney $p$-value,
            yticklabel style={
                align=right,
            },
            ytick={0,0.025,0.05,1},
            % yticklabels={0,0.025,0.05},
        },
        yticklabel style={
            align=right,
            /pgf/number format/fixed,
            /pgf/number format/precision=3,
        },
    ]
    \addplot [
        matrix plot*,
        mesh/rows=10,
        point meta=explicit,
    ] table [
        col sep=comma,
        x=frequency,
        y=n,
        meta expr=\thisrow{mw_p_value}
    ] {measurements/latencies_sig.csv};
    \end{axis}

\end{tikzpicture}
    \caption[Mann-Whitney U Test: Significance Results for Measured Latencies]{Mann-Whitney U Test Significance Results for Measured Latencies (Results below 0.05 are Considered Statistically Significant)} \label{app:mwtest}
\end{figure}


\begin{figure}[h]
    \centering
    \begin{tikzpicture}
    \newcommand{\cpuPlotneoj}[2]{
        \addplot+[
            mark=\symbolneoj,
            mark options={color=#2},
            color=#2,
            solid,
            thick
        ] table [
            x=frequency,
            y expr=\thisrow{#1/neo4j/cpu_median}*100/6,
        ] {\datatable};
    }
    \newcommand{\cpuPlotinfluxdb}[2]{
        \addplot+[
            mark=\symbolinfluxDB,
            mark options={color=#2},
            color=#2,
            solid,
            semithick
        ] table [
            x=frequency,
            y expr=\thisrow{#1/influxdb/cpu_median}*100/6,
        ] {\datatable};
    }
    \newcommand{\cpuPlotdatamgmt}[2]{
        \addplot+[
            mark=\symbolShadowCaster,
            mark options={color=#2},
            color=#2,
            solid,
            thick
        ] table [
            x=frequency,
            y expr=\thisrow{#1/datamgmt/cpu_median}*100/6,
        ] {\datatable};
    }

    \begin{axis}[
        name=second,
        xlabel={Frequency [Hz]},
        ylabel=\empty,
        /tikz/every even column/.append style={column sep=0.4cm},
        width=14.5cm,
        height=4.5cm,
        error bars/y dir=both,
        error bars/y explicit,
        xmin=0,
        xmax=100,
        xtick pos=bottom,
        ytick pos=left,
        ytick={0,.25,.5},
        ymin=0,
        ymax=.625,
        % ytick distance=.25,
        grid=both,
        tick align=center,
        minor y tick num=1,
        minor grid style={gray!25},
    ]

        \pgfplotstableread[col sep=comma]{measurements/cpu_usage_trace_reformatted.csv}\datatable
        \cpuPlotinfluxdb{0}{tubaf-energy}
        \cpuPlotinfluxdb{10}{tubaf-geo}
        \cpuPlotinfluxdb{20}{tubaf-primary-dark}
        \cpuPlotinfluxdb{30}{tubaf-material}
        \cpuPlotinfluxdb{40}{tubaf-environment}
        \cpuPlotdatamgmt{0}{tubaf-energy}
        \cpuPlotdatamgmt{10}{tubaf-geo}
        \cpuPlotdatamgmt{20}{tubaf-primary-dark}
        \cpuPlotdatamgmt{30}{tubaf-material}
        \cpuPlotdatamgmt{40}{tubaf-environment}
    \end{axis}
    \begin{axis}[
        at={($ (second.north west) + (0,2) $)},
        anchor=south west,
        xlabel=\empty,
        ylabel=\empty,
        legend style={at={(.015,.98)}, anchor=north west, legend columns=2, font=\footnotesize},
        /tikz/every even column/.append style={column sep=0.4cm},
        width=14.5cm,
        height=8cm,
        error bars/y dir=both,
        error bars/y explicit,
        xticklabels=\empty,
        xtick pos=bottom,
        xtick distance=10,
        xmin=0,
        xmax=100,
        ytick pos=left,
        ytick distance=.25,
        ymin=0,
        ymax=1.6666,
        grid=both,
        tick align=center,
        minor y tick num=1,
        minor grid style={gray!25},
    ]

        \pgfplotstableread[col sep=comma]{measurements/cpu_usage_trace_reformatted.csv}\datatable
        \addlegendimage{color=tubaf-energy, thick}
        \addlegendentry{n=~0}
            \addlegendimage{only marks, mark=\symbolneoj,color=tubaf-black}
            \addlegendentry{neo4j}
        \addlegendimage{color=tubaf-geo, thick}
        \addlegendentry{n=10}
            \addlegendimage{only marks, mark=\symbolinfluxDB,color=tubaf-black}
            \addlegendentry{influxDB}
        \addlegendimage{color=tubaf-primary-dark, thick}
        \addlegendentry{n=20}
            \addlegendimage{only marks, mark=\symbolShadowCaster,color=tubaf-black}
            \addlegendentry{Shadow Caster}
        \addlegendimage{color=tubaf-material, thick}
        \addlegendentry{n=30}
            \addlegendimage{white}
            \addlegendentry{}
        \addlegendimage{color=tubaf-environment, thick}
        \addlegendentry{n=40}


        \cpuPlotneoj{0}{tubaf-energy}
        \cpuPlotneoj{10}{tubaf-geo}
        \cpuPlotneoj{20}{tubaf-primary-dark}
        \cpuPlotneoj{30}{tubaf-material}
        \cpuPlotneoj{40}{tubaf-environment}

    \end{axis}

    \node [above, rotate=90, yshift=1cm] at (second.north west) {CPU Utilisation [\%]};
\end{tikzpicture}
    \caption{CPU Utilisation of Shadow Caster and Shadow Instance} \label{app:cpuUsageStreamDetailed}
\end{figure}

\begin{figure}[h]
    \centering
    \begin{tikzpicture}
    \newcommand{\ramPlotShadowCaster}[2]{
        \addplot+[
            mark=\symbolShadowCaster,
            mark options={color=#2},
            color=#2,
            solid,
            thick
        ] table [
            x=frequency,
            y expr=\thisrow{#1/stream/proportionalResident_median}*1e-6,
            % y error expr=\thisrow{#1/stream/proportionalResident_std}*1e-6,
        ] {\datatable};
    }
    \newcommand{\ramPlotShadowInstance}[2]{
        \addplot+[
            mark=\symbolShadowInstance,
            mark options={color=#2},
            color=#2,
            solid,
            thick
        ] table [
            x=frequency,
            y expr=\thisrow{#1/db/memrss_median}*1e-6,
        ] {\datatable};
    }
    \newcommand{\ramPlotRos}[2]{
        \addplot+[
            mark=\symbolROSSystem,
            mark options={color=#2},
            color=#2,
            solid,
            thick
        ] table [
            x=frequency,
            y expr=\thisrow{#1/ros2/proportionalResident_median}*1e-6,
        ] {\datatable};
    }

    \begin{axis}[
        name=third,
        xlabel={Frequency [Hz]},
        ylabel=\empty,
        legend style={at={(.99,.09)}, anchor=east, legend columns=5, font=\footnotesize},
        /tikz/every even column/.append style={column sep=0.4cm},
        width=13cm,
        height=8cm,
        error bars/y dir=both,
        error bars/y explicit,
        xmin=0,
        xmax=100,
        xtick pos=bottom,
        ytick pos=left,
        ytick distance=250,
        minor y tick num=1,
        ymin=0,
        % ymax=100,
        grid=both,
        tick align=center,
        minor grid style={gray!25},
    ]

        \pgfplotstableread[col sep=comma]{measurements/ram_usage_trace_reformatted.csv}\datatable

        \addlegendimage{color=tubaf-energy, thick, mark=\symbolROSSystem}
        \addlegendentry{n=0}
        \addlegendimage{color=tubaf-geo, thick, mark=\symbolROSSystem}
        \addlegendentry{n=10}
        \addlegendimage{color=tubaf-primary-dark, thick, mark=\symbolROSSystem}
        \addlegendentry{n=20}
        \addlegendimage{color=tubaf-material, thick, mark=\symbolROSSystem}
        \addlegendentry{n=30}
        \addlegendimage{color=tubaf-environment, thick, mark=\symbolROSSystem}
        \addlegendentry{n=40}

        \ramPlotRos{0}{tubaf-energy}
        \ramPlotRos{10}{tubaf-geo}
        \ramPlotRos{20}{tubaf-primary-dark}
        \ramPlotRos{30}{tubaf-material}
        \ramPlotRos{40}{tubaf-environment}
    \end{axis}
    \begin{axis}[
        name=second,
        at={($ (third.north west) + (0,2) $)},
        anchor=south west,
        xlabel=\empty,
        ylabel={RAM Utilisation [MB]},
        legend style={at={(.99,.02)}, anchor=south east, legend columns=5, font=\footnotesize},
        /tikz/every even column/.append style={column sep=0.4cm},
        width=13cm,
        height=8cm,
        error bars/y dir=both,
        error bars/y explicit,
        xmin=0,
        xmax=100,
        xticklabels=\empty,
        xtick pos=bottom,
        ytick pos=left,
        ytick distance=250,
        minor y tick num=1,
        % ytick={0,1,2,3,4,5},
        ymin=500,
        % ymax=100,
        grid=both,
        tick align=center,
        minor grid style={gray!25},
    ]

        \pgfplotstableread[col sep=comma]{measurements/ram_usage_trace_reformatted.csv}\datatable

        \addlegendimage{color=tubaf-energy, thick, mark=\symbolShadowInstance}
        \addlegendentry{n=0}
        \addlegendimage{color=tubaf-geo, thick, mark=\symbolShadowInstance}
        \addlegendentry{n=10}
        \addlegendimage{color=tubaf-primary-dark, thick, mark=\symbolShadowInstance}
        \addlegendentry{n=20}
        \addlegendimage{color=tubaf-material, thick, mark=\symbolShadowInstance}
        \addlegendentry{n=30}
        \addlegendimage{color=tubaf-environment, thick, mark=\symbolShadowInstance}
        \addlegendentry{n=40}

        \ramPlotShadowInstance{0}{tubaf-energy}
        \ramPlotShadowInstance{10}{tubaf-geo}
        \ramPlotShadowInstance{20}{tubaf-primary-dark}
        \ramPlotShadowInstance{30}{tubaf-material}
        \ramPlotShadowInstance{40}{tubaf-environment}
    \end{axis}
    \begin{axis}[
        name=first,
        at={($ (second.north west) + (0,2) $)},
        anchor=south west,
        xlabel=\empty,
        ylabel=\empty,
        legend style={at={(.99,.02)}, anchor=south east, legend columns=5, font=\footnotesize},
        /tikz/every even column/.append style={column sep=0.4cm},
        width=13cm,
        height=8cm,
        error bars/y dir=both,
        error bars/y explicit,
        xmin=0,
        xmax=100,
        xticklabels=\empty,
        xtick pos=bottom,
        ytick pos=left,
        ytick={21,23,25},
        minor y tick num=1,
        ymin=21,
        ymax=25,
        grid=both,
        tick align=center,
        minor grid style={gray!25},
    ]

        \pgfplotstableread[col sep=comma]{measurements/ram_usage_trace_reformatted.csv}\datatable

        \addlegendimage{color=tubaf-energy, thick, mark=\symbolShadowCaster}
        \addlegendentry{n=0}
        \addlegendimage{color=tubaf-geo, thick, mark=\symbolShadowCaster}
        \addlegendentry{n=10}
        \addlegendimage{color=tubaf-primary-dark, thick, mark=\symbolShadowCaster}
        \addlegendentry{n=20}
        \addlegendimage{color=tubaf-material, thick, mark=\symbolShadowCaster}
        \addlegendentry{n=30}
        \addlegendimage{color=tubaf-environment, thick, mark=\symbolShadowCaster}
        \addlegendentry{n=40}

        \ramPlotShadowCaster{0}{tubaf-energy}
        \ramPlotShadowCaster{10}{tubaf-geo}
        \ramPlotShadowCaster{20}{tubaf-primary-dark}
        \ramPlotShadowCaster{30}{tubaf-material}
        \ramPlotShadowCaster{40}{tubaf-environment}
    \end{axis}


    \node [below, xshift=5, rotate=90] at (first.east) {Shadow Caster};
    \node [below, xshift=5, rotate=90] at (second.east) {Shadow Instance};
    \node [below, xshift=5, rotate=90] at (third.east) {ROS2 System};
\end{tikzpicture}
    \caption{Detailed RAM Utilisation of Shadow Caster, Shadow Instance, and ROS2 system} \label{app:ramUsageDetailed}
\end{figure}


\begin{figure}[h]
    \centering
    \begin{tikzpicture}
    \begin{axis}[
        % ymode=logd,
        xlabel={Frequency [Hz]},
        ylabel={Relative CPU Utilisation ($\frac{\text{Shadow Implementation}*100}{\text{ROS2 System}}$) [\%]},
        legend style={at={(.98,.5)}, anchor=east, font=\footnotesize},
        /tikz/every even column/.append style={column sep=0.4cm},
        width=14cm,
        height=10cm,
        error bars/y dir=both,
        error bars/y explicit,
        xtick pos=bottom,
        xmin=10,
        xmax=100,
        ytick pos=left,
        ymin=0.1,
        ymax=200,
        ytick distance=50,
        grid=both,
        tick align=center,
        minor y tick num=1,
        minor grid style={gray!25},
    ]

        \pgfplotstableread[col sep=comma]{measurements/cpu_usage_trace_reformatted.csv}\datatable

        \pgfplotstableread[col sep=comma]{measurements/cpu_usage_trace_reformatted.csv}\datatable
 
        \addlegendimage{color=tubaf-geo, thick}
        \addlegendentry{n=10}
        \addlegendimage{color=tubaf-primary-dark, thick}
        \addlegendentry{n=20}
        \addlegendimage{color=tubaf-material, thick}
        \addlegendentry{n=30}
        \addlegendimage{color=tubaf-environment, thick}
        \addlegendentry{n=40}


        \newcommand{\cpuPlotCaster}[2]{
            \addplot+[
                mark=none,
                color=#2,
                solid,
                thick
            ] table [
                x=frequency,
                y expr=(\thisrow{#1/stream/cpu_median}/\thisrow{#1/ros2/cpu_median})*100,
            ] {\datatable};
        }
        \cpuPlotCaster{10}{tubaf-geo}
        \cpuPlotCaster{20}{tubaf-primary-dark}
        \cpuPlotCaster{30}{tubaf-material}
        \cpuPlotCaster{40}{tubaf-environment}

    \end{axis}
\end{tikzpicture}
    \caption[Relative CPU Utilisation of Shadow Implementation]{Relative CPU Utilisation of Shadow Implementation} \label{app:cpuUsageStreamRelative}
\end{figure}

\begin{figure}[h]
    \centering
    \begin{tikzpicture}
    \begin{axis}[
        name=third,
        xlabel={Number of Publisher-Subscriber Pairs (n)},
        ylabel={Relative RAM Utilisation ($\frac{\text{Entity}*100}{\text{ROS2 System}}$) [\%]},
        legend style={at={(.98,.5)}, anchor= east, font=\footnotesize},
        /tikz/every even column/.append style={column sep=0.4cm},
        width=14cm,
        height=10cm,
        error bars/y dir=both,
        error bars/y explicit,
        xmin=10,
        xmax=40,
        xtick distance=10,
        xtick pos=bottom,
        ytick pos=left,
        % ytick distance=500,
        minor x tick num=0,
        minor y tick num=1,
        ymin=0,
        grid=both,
        tick align=center,
        minor grid style={gray!25},
    ]

        \pgfplotstableread[col sep=comma]{measurements/ram_usage_aggregated_trace.csv}\datatable

    \addplot+[
        mark=none,
        mark options={color=tubaf-geo},
        color=tubaf-geo,
        solid,
        very thick
    ] table [
        x=num_pairs,
        y expr=((\thisrow{influxdb/memrss_mean}+\thisrow{neo4j/memrss_mean})/(\thisrow{ros2talker/proportionalResident_mean}+\thisrow{ros2listener/proportionalResident_mean}))*100,
    ] {\datatable};
    \addlegendentry{Shadow Instance}

    \addplot+[
        mark=none,
        mark options={color=tubaf-material},
        color=tubaf-material,
        solid,
        very thick
    ] table [
        x=num_pairs,
        y expr=((\thisrow{datamgmt/proportionalResident_mean})/(\thisrow{ros2talker/proportionalResident_mean}+\thisrow{ros2listener/proportionalResident_mean}))*100,
    ] {\datatable};
    \addlegendentry{Shadow Caster}

    \end{axis}

\end{tikzpicture}
    \caption{Relative RAM Utilisation of Shadow Caster and Shadow Instance} \label{app:ramUsageStreamRelative}
\end{figure}

\begin{figure}[h]
    \centering
    \begin{tikzpicture}
    \pgfplotstableread[col sep=comma]{measurements/cpu_usage_rotationLive.csv}\datatable
    \begin{axis}[
        name=first,
        ylabel={CPU Utilisation [\%]},
        xlabel={Time running [min]},
        legend style={at={(.02,.575)}, anchor=west, font=\footnotesize},
        boxplot/draw direction=y,
        xmajorgrids=true,
        xticklabel style={font=\small},
        % xtick={1,2},
        xtick pos=bottom,
        % xticklabels={LTTng-Live,Rotation},
        xmin=0,
        xmax=12,
        ymin=0,
        ymax=8.5,
        ytick distance=2.5,
        minor y tick num=1,
        yticklabel style={font=\small},
        width=9cm,
        height=10cm,
        line width=.5pt,
        ytick pos=left,
        tick align=center,
        grid=both,
    ]

    \addplot+[ 
        mark=none,
        color=tubaf-energy,
        solid,
        thick
    ] table [
        x expr=\thisrow{t}/60,
        y expr=\thisrow{structuralLive}*100/6,
    ] {\datatable};
    \addplot+[
        mark=none,
        color=tubaf-primary,
        solid,
        thick
    ] table [
        x expr=\thisrow{t}/60,
        y expr=\thisrow{sessiond}*100/6,
    ] {\datatable};

    \addlegendimage{color=tubaf-energy, thick}
    \addlegendentry{LTTng-Live}
    \addlegendimage{color=tubaf-primary, thick}
    \addlegendentry{Rotation}
    \end{axis}
    \begin{axis}[
        at={(first.north east)},
        anchor=north west,
        ylabel=\empty,
        boxplot/draw direction=y,
        xmajorgrids=true,
        xtick={1,2},
        xticklabel=\empty,
        xmajorticks=false,
        ymin=0,
        ymax=8.5,
        yticklabel=\empty,
        ytick distance=2.5,
        minor y tick num=1,
        width=4cm,
        height=10cm,
        line width=.5pt,
        ytick pos=left,
        tick align=center,
        grid=both,
    ]

    \addplot+[
        name=second,
        at={(first.north east)},
        anchor=north west,
        boxplot,
        draw=tubaf-energy,
        thick,
    ] table[
        col sep=comma,
        y expr=\thisrow{structuralLive}*100/6,
    ] {measurements/cpu_usage_rotationLive.csv} [right, tubaf-black]
        node at
        (boxplot box cs:\boxplotvalue{median},1)
        {\footnotesize\pgfmathprintnumber{\boxplotvalue{median}}};

    \addplot+[
        boxplot,
        draw=tubaf-primary,
        thick,
    ] table[
        col sep=comma,
        y expr=\thisrow{sessiond}*100/6,
    ] {measurements/cpu_usage_rotationLive.csv} [left, tubaf-black]
        node at
        ($(boxplot box cs:\boxplotvalue{median},1) + (-.2,.1)$)
        % {\footnotesize\pgfmathprintnumber{\boxplotvalue{median}}};
        {\footnotesize\pgfmathprintnumber[fixed,precision=4]{0.0833}};

    \end{axis}
\end{tikzpicture}
    \caption{CPU Utilisation for Differing Tracer Implementations} \label{app:eventProbLatencyCPU}
\end{figure}

\begin{figure}[h]
    \centering
    \begin{tikzpicture}
    \begin{axis}[
        name=second,
        xlabel={Relative Time [min]},
        ylabel=\empty,
        tick align=center,
        legend style={
            at={(.99,.05)},
            anchor=south east,
            legend columns=3,
            draw=black,
            font=\footnotesize
        },
        xtick distance=1,
        % minor ytick={0.5,1.5,...,5},
        % ytick distance=1,
        xmin=0, xmax=5.9,
        ymin=17, ymax=20,
        ytick={17, 18, 19},
        grid=both,
        height=4cm,
        width=14cm,
        xtick pos=bottom,
        ytick pos=left,
    ]

        \addplot+ [mark=none, color=tubaf-geo, thick] table [
            x expr=(\thisrow{timeRelative}/60)-6.5,
            y expr=(\thisrow{datamgmt/proportionalResident}+\thisrow{continuous/proportionalResident}+\thisrow{structural/proportionalResident})*1e-6,
            col sep=comma,
            mark=none
        ] {measurements/totaling_overhead_ram.csv};

        \legend{
            {Shadow Caster},
        }

        \path (2.74,17) coordinate (bottom);

    \end{axis}
    \begin{axis}[
        at={($(second.north west) + (0,2)$)},
        anchor=south west,
        xlabel=\empty,
        ylabel=\empty,
        tick align=center,
        legend style={
            at={(.99,.5)},
            anchor=south east,
            draw=black,
            font=\footnotesize
        },
        xtick distance=1,
        % minor ytick={0.5,1.5,...,5},
        % ytick distance=1,
        xticklabels=\empty,
        xmin=0, xmax=5.9,
        ymin=0, ymax=1050,
        grid=both,
        height=7cm,
        width=14cm,
        xtick pos=bottom,
        ytick pos=left,
    ]

        \addplot+ [mark=none, color=tubaf-material, thick] table [
            x expr=(\thisrow{timeRelative}/60)-6.5,
            y expr=\thisrow{neo4j/memrss}*1e-6,
            col sep=comma,
            mark=none,
        ] {measurements/totaling_overhead_ram.csv};
        \addplot+ [mark=none, color=tubaf-environment, thick] table [
            x expr=(\thisrow{timeRelative}/60)-6.5,
            y expr=\thisrow{influxdb/memrss}*1e-6,
            col sep=comma,
            mark=none,
        ] {measurements/totaling_overhead_ram.csv};


        \legend{
            {neo4j},
            {influxDB},
        }


        \path (2.74,1050) coordinate (top);
    \end{axis}


    \draw [tubaf-primary-dark, semithick, dashed] (bottom) -- (top) node [midway, above left] {started request};

    \node [above, rotate=90, yshift=1.5cm] at (second.north west) {RAM Utilisation [MB]};

\end{tikzpicture}
    \caption{Influence of Requests on the RAM Utilisation of the Shadow Implementation} \label{app:aggregationOverheadRAM}
\end{figure}


\chapter{Real World Scenario}
    \begin{table}[h]
        \centering
        \caption{List of Active Nodes in Stage 3 of the Real-World Scenario}
        \begin{tblr}{ll}
            \toprule
            /emergency\_stop    &   /lynx/system\_monitor       \\
            /fixposition\_driver\_ros2  &   /twist\_mux \\
            /foxglove\_bridge   &   /velocity\_smoother \\
            /foxglove\_bridge\_component\_manager   &   /wrapper/fixposition\_wrapper\_node \\
            /joy\_teleop/joy\_node  &   /wrapper/livox\_front\_wrapper\_node \\
            /joy\_teleop/teleop\_twist\_joy\_node   &   /wrapper/sick\_front\_wrapper\_node \\
            /launch\_ros\_1 &   /wrapper/sick\_rear\_wrapper\_node \\
            /launch\_ros\_143238    &   /zed2i\_front/transform\_listener\_impl\_... \\
            /launch\_ros\_82636 &   /zed2i\_front/zed2i\_front\_state\_publisher \\
            /lifecycle\_manager\_rosee  &   /zed2i\_front/zed\_container \\
            /lynx/battery\_driver   &   /zed2i\_front/zed\_node\_front \\
            /lynx/controller\_manager   &   /zed2i\_left/transform\_listener\_impl\_... \\
            /lynx/drive\_controller &   /zed2i\_left/zed2i\_left\_state\_publisher \\
            /lynx/hardware\_controller  &   /zed2i\_left/zed\_container \\
            /lynx/imu\_broadcaster  &   /zed2i\_left/zed\_node\_left \\
            /lynx/joint\_state\_broadcaster &   /zed2i\_rear/transform\_listener\_impl\_... \\
            /lynx/joy2twist\_node   &   /zed2i\_rear/zed2i\_rear\_state\_publisher \\
            /lynx/joy\_node &   /zed2i\_rear/zed\_container \\
            /lynx/lights\_container &   /zed2i\_rear/zed\_node\_rear \\
            /lynx/lights\_controller    &   /zed2i\_right/transform\_listener\_impl\_... \\
            /lynx/lights\_driver    &   /zed2i\_right/zed2i\_right\_state\_publisher \\
            /lynx/lights\_manager   &   /zed2i\_right/zed\_container \\
            /lynx/robot\_state\_publisher   &   /zed2i\_right/zed\_node\_right \\
            /lynx/safety\_manager   &   \\
            \bottomrule
        \end{tblr}
    \end{table}

It is imperativ to note, that only those entities are represented within the Graph, that are executed by the monitored host. Extending the Shadow Implementation to multi-host is addressed in Future Work.

\newpage
\section*{remove me}
\section*{remove me}
\section*{remove me}
\begin{figure}
    \caption{Real-World Scenario: Generated ROS2 Graph (Stage 3)} \label{app:graph2}.
\end{figure}
\cleardoublepage